\section{The correlation theorem and congruence admissibility}
Throughout, \(L_0\ge2\) is even,
\(G=(\Z/L_0\Z)^\times\), and \(H_1,H_2\le G\) are fixed.
Define
\[
 \mathcal P_i=\{p\nmid L_0:p\bmod L_0\in H_i\},\quad
 C_i=\langle\mathcal P_i\rangle,\quad
 \rho_i=\frac{|H_i|}{\varphi(L_0)},\quad
 \kappa=2-\rho_1-\rho_2.
\]
An empty product is one. Fix \(U\subset H_1\), \(V\subset H_2\),
and coprime positive \(L_0\)-supported integers \(A,B\).
Write
\[
 R_{U,V}^{A,B}(N;I)=\#\{Ax+By=N:x\in C_1,\ y\in C_2,
      \ x\bmod L_0\in U,\ y\bmod L_0\in V,\ Ax/N\in I\}.
\]
Omit \(I\) when \(I=[0,1]\). Put
\begin{align}
 c_i&=\frac{1}{\Gamma(\rho_i)}
    \lim_{s\downarrow1}(s-1)^{\rho_i}
              \prod_{p\in\mathcal P_i}(1-p^{-s})^{-1},\label{eq:marginal}\\
 \Gamma_N&=\frac{L_0}{|H_1||H_2|}
    \#\{(u,v)\in U\times V:Au+Bv\equiv N\pmod{L_0}\},\label{eq:finite}\\
 \mathcal E&=\{p\nmid L_0:p\notin\mathcal P_1,
                                  \ p\notin\mathcal P_2\},\notag\\
 C_{\mathcal E}&=\prod_{p\in\mathcal E}
               \left(1-\frac{1}{(p-1)^2}\right),\qquad
 J_{\mathcal E}(N)=\prod_{\substack{p\mid N\\p\in\mathcal E}}
                              \frac{p-1}{p-2}.\label{eq:singular}
\end{align}
The infinite product in \eqref{eq:singular} converges absolutely and
is positive, since \(2\mid L_0\). The constants \(c_i\) are positive.
For an explicit convergent expression, let \(r_i(p)>1\) denote the
order of \(pH_i\) in \(G/H_i\) at a forbidden good prime. Then
\begin{equation}\label{eq:marginal-product}
 c_i=\frac{(\varphi(L_0)/L_0)^{\rho_i}}{\Gamma(\rho_i)}
 \prod_{\substack{\eta\in H_i^\perp\\\eta\ne1}}L(1,\eta)^{\rho_i}
 \prod_{\substack{p\nmid L_0\\p\notin\mathcal P_i}}
                     (1-p^{-r_i(p)})^{1/r_i(p)}.
\end{equation}
Characters here are literal modulo \(L_0\). Continuation along the
positive real axis from \(s>1\) fixes the powers; conjugate factors
are paired. The Euler identity proving this formula is given in
Section~\ref{sec:model}.

\begin{theorem}\label{thm:general}
With all these data fixed, let
\(\delta=1/(8192\varphi(L_0))\). Uniformly in \(N\) and every interval
\(I\subset[0,1]\),
\begin{align*}
 R_{U,V}^{A,B}(N;I)
 &=c_1c_2\Gamma_N C_{\mathcal E}|I|
       \frac{NJ_{\mathcal E}(N)}{AB(\log N)^\kappa}\\
 &\quad+O\left(\frac{NJ_{\mathcal E}(N)}
                         {AB(\log N)^{\kappa+\delta}}\right).
\end{align*}
If \(\Gamma_N=0\), the count is identically zero.
For \(\Gamma_N>0\) and \(I=[0,1]\), the error is a relative
\(O((\log N)^{-\delta})\). More generally it is relative \(o(1)\)
when \(|I|(\log N)^\delta\to\infty\).
\end{theorem}

The implicit constant and threshold may depend on all fixed data.
No numerical threshold, uniformity in a growing modulus, or
power-short interval assertion is included. For noncoprime fixed
coefficients, first divide by their gcd when it divides \(N\).
The even-modulus convention makes the prime two a bad prime throughout
the analytic theorem. Passing from an odd modulus to its even multiple
requires retaining this convention; freely admitting the prime two is
a separate valuation-profile operation, not a change of notation.

\begin{theorem}\label{thm:c12}
Let \(\mathcal E_{12}=\{p\nmid6:p\not\equiv1\pmod{12}\}\), and set
\[
 A_{12}=\frac{4\pi}{\Gamma(1/4)^2}
       \sqrt{\frac{\log(2+\sqrt3)}{6}}
       \prod_{p\in\mathcal E_{12}}
                        \left(1-\frac{2}{p(p-1)}\right),\qquad
 J_{12}(N)=\prod_{\substack{p\mid N\\p\in\mathcal E_{12}}}
                                      \frac{p-1}{p-2}.
\]
For all sufficiently large \(N\equiv2\pmod{12}\),
\[
 R_{12}(N)=A_{12}J_{12}(N)\frac{N}{(\log N)^{3/2}}
                 \left(1+O((\log N)^{-1/32768})\right).
\]
All other residue classes have zero count.
\end{theorem}
\begin{proof}[Deduction from Theorem~\ref{thm:general}]
Take \(L_0=12\), \(H_1=H_2=U=V=\{1\}\), \(A=B=1\).
Then \(\rho_i=1/4\) and \(\Gamma_N=12\) on the indicated class.
The four characters give the exact identity
\begin{align*}
 F_{12}(s)^4={}&\zeta(s)L(s,\chi_{-4})L(s,\chi_{-3})L(s,\chi_{12})\,
        (1-2^{-2s})(1-3^{-2s})
        \prod_{p\in\mathcal E_{12}}(1-p^{-2s})^2.
\end{align*}
Here the three nonprincipal characters displayed are primitive.
Use \(L(1,\chi_{-4})=\pi/4\),
\(L(1,\chi_{-3})=\pi/(3\sqrt3)\), and
\[
 L(1,\chi_{12})=\int_0^1\frac{1-x^4}{1+x^6}\,dx
 =\left[\frac{1}{2\sqrt3}
 \log\frac{x^2+\sqrt3x+1}{x^2-\sqrt3x+1}\right]_0^1
 =\frac{\log(2+\sqrt3)}{\sqrt3}.
\]
Consequently
\[
 c_{12}=\frac{1}{\Gamma(1/4)}
 \left(\frac{\pi^2\log(2+\sqrt3)}{54}
       \prod_{p\in\mathcal E_{12}}(1-p^{-2})^2\right)^{1/4}.
\]
In \(12c_{12}^2C_{\mathcal E_{12}}\), combine the factors by
\((1-p^{-2})(1-(p-1)^{-2})=1-2/(p(p-1))\).
This is exactly \(A_{12}\). The real density integrates to one;
there is no extra beta-integral factor.
\end{proof}

The uniformity in the target also permits a coprimality restriction.
Only the primes allowed in the summands can divide their gcd.
Removing them changes the leading term by the following finite product.
\begin{corollary}\label{cor:coprime-c12}
Put
\[
 R_{\mathrm{cop}}(N)=\#\{(a,b)\in C_{12}^2:a+b=N,\ (a,b)=1\},
 \qquad
 P_{12}(N)=\prod_{\substack{p\mid N\\p\equiv1\ (12)}}(1-p^{-1}).
\]
Uniformly over sufficiently large \(N\equiv2\pmod{12}\),
\[
 R_{\mathrm{cop}}(N)
 =A_{12}J_{12}(N)P_{12}(N)\frac{N}{(\log N)^{3/2}}
       \left(1+O((\log N)^{-1/65536})\right).
\]
Every other residue class has zero count. In particular, every
sufficiently large admissible target has coprime summands.
\end{corollary}
\begin{proof}
M\"obius inversion gives the exact identity
\begin{equation}\label{eq:coprime-inversion}
 R_{\mathrm{cop}}(N)
 =\sum_{\substack{d\mid N\\d\in C_{12}}}\mu(d)R_{12}(N/d).
\end{equation}
Indeed, division by any common divisor preserves both prime supports,
and every common divisor lies in \(C_{12}\). Moreover \(d\equiv1\pmod{12}\),
so \(N/d\equiv2\pmod{12}\), and
\(J_{12}(N/d)=J_{12}(N)\).

Write \(L=\log N\), \(\ell=\log\log(3N)\),
\(P=P_{12}(N)\), and \(\delta=1/32768\).
We use the elementary divisor estimates
\begin{align}
 S_0&=\sum_{\substack{d\mid N\\d\in C_{12}}}\frac{|\mu(d)|}{d}
       =\prod_{\substack{p\mid N\\p\equiv1\ (12)}}(1+p^{-1})
       \ll\ell,\qquad P^{-1}\ll\ell,\label{eq:coprime-products}\\
 S_1&=\sum_{\substack{d\mid N\\d\in C_{12}}}
                          \frac{|\mu(d)|\log d}{d}
       =S_0\sum_{\substack{p\mid N\\p\equiv1\ (12)}}
                                     \frac{\log p}{p+1}
       \ll\ell^3.\label{eq:coprime-logmoment}
\end{align}
For \eqref{eq:coprime-products}, split the prime products at \(L\).
Mertens' product bound controls primes at most \(L\) by \(O(\log L)\);
the remaining factors have bounded product, since
\(\sum_{p\mid N,\ p>L}p^{-1}\le1/\log L\).
For \eqref{eq:coprime-logmoment}, the primes at most \(L\) contribute
at most \(\sum_{n\le L}(\log n)/n\ll(\log L)^2\), and the others
contribute at most \(L^{-1}\sum_{p\mid N}\log p\le1\).

Split \eqref{eq:coprime-inversion} at \(d=\sqrt N\).
The bound \(R_{12}(M)\le M\) shows that its large-divisor part is
\(O(\sqrt N\,\tau(N))\), where \(\tau\) is the divisor function.
For \(d\le\sqrt N\), Theorem~\ref{thm:c12} applies uniformly, and
\[
 (L-\log d)^{-3/2}
       =L^{-3/2}+O(L^{-5/2}\log d).
\]
Consequently, using \(\sum_{d\mid N,\ d\in C_{12}}\mu(d)/d=P\)
and extending the main sum back to all divisors,
\begin{align*}
 R_{\mathrm{cop}}(N)
 ={}&A_{12}J_{12}(N)\frac{N}{L^{3/2}}
       \left(P+O(L^{-\delta}S_0+L^{-1}S_1
                            +N^{-1/2}\tau(N))\right)\\
 &\quad+O(\sqrt N\,\tau(N)).
\end{align*}
After division by the main term, the error is
\[
 O\!\left(L^{-\delta}\ell^2+L^{-1}\ell^4
                  +N^{-1/2}\tau(N)L^{3/2}\ell\right).
\]
Here \(J_{12}(N)\ge1\) and \(\tau(N)=N^{o(1)}\).
The displayed error is \(O(L^{-\delta/2})\), proving the claim.
\end{proof}

\begin{remark}[Comparison with the quartic-norm target weight]
The weight in \cite[Theorem 2]{iwaniecpomykala} is
\(W_{12}(N)=\prod_{p\mid N}(1-\varepsilon_p/p)\), with
\(\varepsilon_p=1\) on allowed primes and \(-1\) otherwise.
On the admissible class,
\[
 \frac{J_{12}(N)P_{12}(N)}{W_{12}(N)}
 =\frac23\prod_{\substack{p\mid N\\p\equiv5,7,11\ (12)}}
                 \frac{p(p-1)}{(p-2)(p+1)}\asymp1
\]
uniformly in \(N\), since the factors are \(1+O(p^{-2})\).
The factor \(2/3\) comes from \(2\mid N\); the prime three does
not divide an admissible target. Thus the coprime asymptotic has
the same order as the earlier conditional coprime bounds.
\end{remark}

We next clarify the congruence criterion without analytic inputs.
Write \(C_H(U)=\{n\in C_H:n\bmod L_0\in U\}\).
\begin{lemma}\label{lem:finite-images}
A residue \(r\bmod M\) is represented by \(C_H(U)\) if and only if
its reduction modulo \(\gcd(M,L_0)\) is that of some \(h\in U\),
and every forbidden prime \(p\mid M\), \(p\nmid L_0\), satisfies
\(p\nmid r\).
\end{lemma}
\begin{proof}
Necessity follows from prime support. Choose such an \(h\), and for
each allowed \(p^k\Vert M\) put \(e_p=\min(v_p(r),k)\).
Set \(D=\prod p^{e_p}\), so \(D\in C_H\) and \((D,L_0)=1\).
We seek a prime \(\ell\equiv hD^{-1}\pmod{L_0}\).
Where \(e_p<k\), impose
\[
 \ell\equiv(r/p^{e_p})(D/p^{e_p})^{-1}\pmod{p^{k-e_p}}.
\]
Where \(e_p=k\), impose \(\ell\equiv1\pmod p\).
At primes of \(M\) not dividing \(D\), impose the unit condition
\(\ell\equiv rD^{-1}\pmod{p^k}\). These congruences are compatible
with the \(L_0\)-condition precisely by the stated hypotheses.
They give a reduced CRT class. Dirichlet's theorem supplies a prime
\(\ell\) outside the finite exceptional set; then \(D\ell\) has
the required residue and state.
\end{proof}

\begin{theorem}\label{thm:local}
The condition \(\Gamma_N>0\) is equivalent to solvability of
\(Ax+By\equiv N\pmod M\) with actual
\(x\in C_{H_1}(U),y\in C_{H_2}(V)\), for every \(M\ge1\).
For all sufficiently large \(N\), it is also equivalent to positive
additive representability.
\end{theorem}
\begin{proof}
Taking \(M=L_0\) proves necessity. For sufficiency choose a root
counted by \(\Gamma_N\) and work modulo \(K=\lcm(M,L_0)\).
At a prime dividing \(L_0\), at least one coefficient is a unit;
adjust its variable to lift the equation while preserving the fixed
unit states. At a prime outside \(L_0\), both coefficients are units
and \(p\ge3\). Of the \(p\) solutions modulo \(p\), requiring a
unit on each forbidden side removes at most two. A surviving pair
lifts to every higher power. CRT and Lemma~\ref{lem:finite-images}
realize the required two residues by actual semigroup values.
Finally Theorem~\ref{thm:general} is positive for all sufficiently
large targets with \(\Gamma_N>0\).
\end{proof}

\begin{remark}[A joint congruence obstruction]\label{rem:coupled}
Take \(L_0=12\), \(H=\{1,11\}=\ker\chi_{12}\).
Then \(H+H=\{0,2,10\}\pmod{12}\), so \(N\equiv4\pmod{12}\)
is never a sum. Nevertheless \(H\) projects onto the full unit groups
modulo four and three. The separate closures are \(\Z_2^\times\)
and \(\Z_3^\times\); at other primes they are \(\Z_p\) or
\(\Z_p^\times\), according as that prime is allowed or forbidden.
Two units sum to every element of \(\Z_p\) for odd \(p\), and to
every even element for \(p=2\). Thus these targets pass every
separate \(p\)-adic test. The missing condition is the simultaneous
residue state modulo twelve.
\end{remark}
